\SourceLeaf{75}\emph{Ex. gr.} $3^8=6561\equiv1\,(\operatorname{mod.}32)$.

Quivis enim numerus impar vel sub forma $1+4h$, vel sub hac $-1+4h$ comprehenditur: unde propositio protinus sequitur (theor. art. 86).

Quoniam igitur exponens ad quem quicunque numerus impar secundum modulum $2^n$ pertinet, divisor ipsius $2^{n-2}$ esse debet, quivis ad aliquem horum numerum pertinebit $1,2,4,8\mathinner{.\,.\,.\,.}2^{2n-2}$, ad quemnam vero pertineat ita facile diiudicatur. Sit numerus propositus $=4h\pm1$, atque exponens maximae potestatis numeri 2, quae ipsum $h$ metitur, $=m$ (qui etiam $=0$ esse potest, quando scilicet $h$ est impar); tum exponens ad quem numerus propositus pertinet, erit $=2^{n-m-2}$, siquidem $n>m+2$; si autem $n=$ vel $<m+2$, numerus propositus est $\equiv\pm1$ adeoque vel ad exponentem 1 vel ad exponentem 2 pertinebit. Numerum enim formae $\pm1+2^{m+2}k$, (quae huic aequivalet, $4h\pm1$) ad potestatem exponentis $2^{n-m-2}$ elevatum unitati secundum modulum $2^n$ congruum fieri, ad potestatem autem exponentis, qui est inferior numeri 2 potestas, incongruum, ex art. 86 nullo negotio deducitur. Numerus itaque quicunque formae $8k+3$ vel $8k+5$ ad exponentem $2^{n-2}$ pertinebit.

\Article{91}

Hinc patet eo sensu quo supra expressionem accepimus, \emph{radices primitivas} hic non dari, nullos scilicet numeros, quorum periodus omnes numeros modulo minores ad ipsumque primos amplectatur. Attamen facile perspicitur, analogon hic haberi. Invenitur enim, numeri formae $8k+3$ potestatem exponentis imparis semper esse formae $8k+3$, potestatem autem exponentis paris, semper formae $8k+1$; nulla igitur potestas formae $8k+5$ aut $8k+7$ esse potest. Quare quum periodus numeri formae $8k+3$, ex $2^{n-2}$ terminis diversis constet, quorum quisque aut formae $8k+3$ aut huius $8k+1$, neque plures huiusmodi numeri modulo minores dentur quam $2^{n-2}$, manifesto quivis numerus formae $8k+1$ vel $8k+3$ congruus est secundum modulum $2^n$ potestati alicui numeri cujuscunque formae $8k+3$. Simili modo ostendi potest periodum numeri formae $8k+5$ comprehendere omnes numeros formarum $8k+1$ et $8k+5$. Si igitur numerus formae $8k+5$ pro basi assumitur, omnes numeri formae $8k+1$ et $8k+5$, positive, omnesque formae $8k+3$ et $8k+7$, negative sumti, indices reales nanciscentur, et quidem hic indices secundum $2^{n-2}$ congrui pro aequivalentibus sunt habendi. Hoc modo tabula nostra I intelligenda, ubi pro modulis 16, 32 et 64
\SourceLeaf{76}(namque pro modulo $8$ nulla tabula necessaria erit) semper numerum $5$ pro basi accepimus. \emph{Ex. gr.} numero $19$ qui est formae $8n+3$ adeoque \emph{negative} sumendus, respondet pro modulo $64$ index $7$, id quod significat esse $5^7\equiv-19\,(\operatorname{mod.}64)$.

Numeris autem formarum $8n+1$, $8n+5$ negative, atque numeris formarum $8n+3$, $8n+7$ positive acceptis, indices quasi imaginarii tribuendi forent. Quos introducendo calculus indicum ad algorithmum perquam simplicem reduci potest. Sed quoniam, si haec ad omnem rigorem exponere vellemus, nimis longe evagari oporteret, hoc negotium ad aliam occasionem nobis reservamus, quando forsan fusius quantitatum imaginarium theoriam, quae nostro quidem iudicio a nemiue hactenus ad notiones claras est reducta, pertractare suscipiemus. Periti hunc algorithmum facile ipsi eruent: qui minus sunt exercitati, perinde tamen tabula hac uti poterunt, ut ii qui recentiorum commenta de \emph{logarithmis} imaginariis ignorant. logarithmis utuntur, si quidem principia supra stabilita probe tenuerint.

\begin{center}\small\emph{Moduli e pluribus primis compositi.}\end{center}

\Article{92}

Secundum modulum e pluribus primis compositum tantum non omnia quae ad residua potestatum pertinent ex theoria congruentiarum generali deduci possunt; quia vero infra congruentias quascunque secundum modulum e pluribus primis compositum ad congruentias, quarum modulus est primus aut primi potestas, reducere fusius docebimus, non est quod huic rei multum hic immoremur. Observamus tantum, bellissimam proprietatem, quae pro reliquis modulis locum habeat, quod scilicet semper exstent numeri quorum periodus omnes numeros ad modulum primos complectatur, hic deficere, excepto unico casu, quando scilicet modulus est duplum numeri primi, aut potestatis numeri primi.
Si enim modulus $m$ redigitur ad formam $A^aB^bC^c$ etc. designantibus $A,B,C$ etc. numeros primos diversos, praeterea $A^{a-1}(A-1)$ designatur per $\alpha$, $B^{b-1}(B-1)$ per $\beta$ etc. denique $z$ est numerus ad $m$ primus; erit $z^{\alpha}\equiv1\,(\operatorname{mod.}A^a)$. $z^{\beta}\equiv1\,(\operatorname{mod.}B^b)$ etc. Quodsi igitur $\mu$ est minimus numerorum $\alpha,\beta,\gamma$ etc. dividuus communis, erit $z^{\mu}\equiv1$ secundum omnes modulos $A^a,B^b$ etc. adeoque etiam secundum $m$, cui illorum productum est aequale. At excepto casu ubi $m$ est duplum numeri primi aut potestatis numeri primi, numerorum $\alpha,\beta,\gamma$ etc. dividuus communis minimus, ipsorum producto est minor (quoniam numeri $\alpha,\beta,\gamma$ etc. inter se primi esse nequeunt sed certe divisorem $2$ communem habent). Nullius itaque numeri
\SourceLeaf{77}periodus tot terminos comprehendere potest, quot dantur numeri ad modulum primi ipsoque minores, quia horum numerus producto ex $\alpha,\beta,\gamma$ etc. est aequalis. Ita \emph{ex. gr.} pro $m=1001$ cuiusvis numeri ad $m$ primi potestas exponentis $60$ unitati est congrua, quia $60$ est dividuus communis numerorum $6,10,12$. --- Casus autem ubi modulus est duplum numeri primi aut duplum potestatis numeri primi illi ubi est primus aut primi potestas prorsus est similis.

\Article{93}

Scriptorum in quibus alii geometrae de argumento in hac sectione pertractato egerunt, iam passim mentio est facta. Eos tamen qui quaedam fusius, quam nobis brevitas permisit, explicata desiderant, ablegamus imprimis ad sequentes ill. Euleri commentationes, ob perspicuitatem qua vir summus prae omnibus semper excelluit, maxime commendabiles.

\emph{Theoremata circa residua ex divisione potestatum relicta Comm. nov. Petr. T.} VII \emph{p.} 49 sqq.

\emph{Demonstrationes circa residua ex divisione potestatum per numeros primos resultantia.}\par
\hspace*{\parindent}\emph{Ibid. T.} XVIII \emph{p.} 85 sqq.

Adiungi his possunt \emph{Opusculorum analyt. T.} I, \emph{dissertt.} 5 et 8.

\begin{center}\rule{.28\linewidth}{.25pt}\end{center}

\clearpage
\pdfbookmark[0]{Sectio quarta}{section-four}
\SourceLeaf{78}\begin{center}
{\Large SECTIO QUARTA}\par\bigskip
{\small DE}\par\bigskip
{\Large CONGRUENTIIS SECUNDI GRADUS.}
\end{center}

\begin{center}\rule{.09\linewidth}{.25pt}\end{center}

\begin{center}\small\emph{Residua et non-residua quadratica.}\end{center}

\Article{94}

\textsc{Theorema.} \emph{Numero quocunque $m$ pro modulo accepto, ex numeris $0,1,2,3....m-1$, plures quam $\tfrac12m+1$ quando $m$ est par, sive plures quam $\tfrac12m+\tfrac12$, quando $m$ est impar quadrato congrui fieri non possunt.}

\emph{Dem.} Quoniam numerorum congruorum quadrata sunt congrua: quivis numerus, qui ulli quadrato congruus fieri potest, etiam quadrato alicui cuius radix $<m$ congruus erit. Sufficit itaque residua minima quadratorum $0,1,4,9....(m-1)^2$ considerare. At facile perspicitur, esse $(m-1)^2\equiv1$, $(m-2)^2\equiv2^2$, $(m-3)^2\equiv3^2$
etc. Hinc etiam, quando $m$ est par, quadratorum $(\tfrac12m-1)^2$ et $(\tfrac12m+1)^2$, $(\tfrac12m-2)^2$ et $(\tfrac12m+2)^2$ etc. residua minima eadem erunt: quando vero $m$ est impar, quadrata $(\tfrac12m-\tfrac12)^2$ et $(\tfrac12m+\tfrac12)^2$, $(\tfrac12m-\tfrac32)^2$ et $(\tfrac12m+\tfrac32)^2$ etc. erunt congrua. Unde
palam est, alios numeros, quam qui alicui ex quadratis $0,1,4,9....(\tfrac12m)^2$ congrui sint, quadrato congruos fieri non posse, quando $m$ par; quando vero impar, quemvis numerum, qui ulli quadrato sit congruus, alicui ex his $0,1,4,9....(\tfrac12m-\tfrac12)^2$ necessario congruum esse. Quare dabuntur ad summum in priori casu $\tfrac12m+1$ residua minima diversa, in posteriori $\tfrac12m+\tfrac12$. \emph{Q. E. D.}

\emph{Exemplum.} Secundum modulum $13$ quadratorum numerorum $0,1,2,3...6$ residua minima inveniuntur $0,1,4,9,3,12,10$, post haec vero eadem
